\section{Societal, National, and Economic Impact}
\label{sec:impact}

nginx and ZMap sit beneath a large fraction of the open-source software
deployed today, and their defects are therefore inherited by every deployment
built on them.

\paragraph{What depends on nginx.}
nginx sits at the point where untrusted traffic first meets an
organization's software. It terminates TLS, enforces the first layer of
access control, routes requests to the application servers, and mediates
between the internet and the internal identity providers, brokers, and
databases. A defect or a misconfiguration at that point is rarely contained,
because it is a defect in the boundary that every other control assumes to be
intact. Hospitals, payment processors, utilities, and government services all
inherit this exposure, generally without any ability to audit it themselves.

\paragraph{What depends on ZMap.}
ZMap and the tools built on it are how defenders learn what they actually
operate. Attack-surface management, triage after a new advisory, incident
response, and regulatory asset inventory all reduce to the question of what
is running, where, and at what version. When that question is answered incompletely, every downstream decision
inherits the gap, because hosts absent from the inventory are not patched when
an advisory lands. Improving the safety and the
information yield of open-source discovery therefore improves the security of
every organization that relies on it, and it disproportionately helps those
that cannot afford commercial alternatives.

\paragraph{Economic and national significance.}
These components are deployed across the sectors on which the Nation depends,
and the open-source stack as a whole has recently been valued in the
trillions of dollars annually. The maintainer asymmetry described in
Section~\ref{sec:ecosystem} means that public investment in the security of
these specific components generates returns far out of proportion to its cost:
every vulnerability that \sysname finds and drives upstream is removed
simultaneously from every deployment that shares the affected component.

\paragraph{Measures of success.}
The project advances several of the measures of success named in the
solicitation. It establishes new \emph{data sets}, in the reference
infrastructures, their ground-truth knowledge graphs, and the seeded
vulnerability corpus of Section~\ref{sec:evaluation}; \emph{new technologies
and techniques}, in the capture, replication, fidelity-measurement, and
exploitation-chaining methods of Thrusts~1--3, all released as open source;
and new \emph{research infrastructure}, in the twin-construction toolchain
and the public twinCTF infrastructure. Finally, it prepares participants for
careers in \emph{STEM} through the activities of
Section~\ref{sec:broader}.
